\documentclass{article}
\usepackage{graphicx} % Required for inserting images

\usepackage[T1]{fontenc}
\usepackage{enumerate, amsmath, amsfonts, amssymb, amsthm, dsfont, mathrsfs, wasysym, graphics, graphicx, xcolor, url, hyperref, hypcap, xargs, multicol, pdflscape, multirow, hvfloat, array, ae, aecompl, pifont, mathtools, a4wide, float, blkarray, overpic, nicefrac, stmaryrd, anyfontsize, yfonts, array, tabularx}
\usepackage{xargs, bbm, enumerate, paralist, totcount, etoolbox, xfp}
\usepackage[shortlabels, inline]{enumitem}
\usepackage{pdflscape}
\usepackage[noabbrev,capitalise]{cleveref}
\usepackage[normalem]{ulem}
\usepackage{marginnote}

\hypersetup{colorlinks=true, citecolor=darkblue, linkcolor=darkblue}

% math special letters
\newcommand{\R}{\mathbb{R}} % reals
\newcommand{\Q}{\mathbb{Q}} % rationals
\newcommand{\N}{\mathbb{N}} % naturals
\newcommand{\Z}{\mathbb{Z}} % integers
\newcommand{\C}{\mathbb{C}} % complex
% \newcommand{\I}{\mathbb{I}} % set of integers
\newcommand{\HH}{\mathbb{H}} % hyperplane
\newcommand{\K}{k} % field

\newcommand{\f}[1]{{\mathfrak{#1}}} % mathfrak letters
\let\ccc\c
\renewcommand{\c}[1]{{\mathcal{#1}}} % call letters
\renewcommand{\b}[1]{{\boldsymbol{#1}}} % bold letter
\newcommand{\scr}[1]{{\mathscr{#1}}} % curly letter
\newcommand{\go}[1]{{\textgoth{#1}}} % gothic letter

\renewcommand{\emptyset}{\varnothing} % prettier emptyset
\renewcommand{\epsilon}{\varepsilon} % prettier epsilon

% math commands
\newcommand{\set}[2]{\left\{ #1 \;\middle|\; #2 \right\}} % set notation
\newcommand{\bigset}[2]{\big\{ #1 ~;~ #2 \big\}} % big set notation
\newcommand{\Bigset}[2]{\Big\{ #1 ~;~ #2 \Big\}} % Big set notation
\newcommand{\setangle}[2]{\left\langle #1 \;\middle|\; #2 \right\rangle} % set notation
\newcommand{\ssm}{\smallsetminus} % small set minus
\newcommand{\inner}[1]{\left\langle#1\right\rangle}
\newcommand{\dotprod}[2]{\inner{#1,~#2}} % dot product
\newcommand{\bigdotprod}[2]{\big\langle \, #1 \; \big| \; #2 \, \big\rangle} % dot product
\newcommand{\symdif}{\,\triangle\,} % symmetric difference
\newcommand{\eqdef}{\coloneqq}%{\mbox{\,\raisebox{0.2ex}{\scriptsize\ensuremath{\mathrm:}}\ensuremath{=}\,}} % :=
\newcommand{\defeq}{\mbox{~\ensuremath{=}\raisebox{0.2ex}{\scriptsize\ensuremath{\mathrm:}} }} % =:
\renewcommand{\implies}{\Rightarrow} % imply sign

% operators
\DeclareMathOperator{\conv}{conv} % convex hull
\DeclareMathOperator{\vect}{vect} % linear span
\DeclareMathOperator{\cone}{cone} % cone hull
\DeclareMathOperator{\vol}{vol} % volume


% others
\newcommand{\ie}{\textit{i.e.}~} % id est
\newcommand{\eg}{\textit{e.g.}~} % exempli gratia
\newcommand{\Eg}{\textit{E.g.}~} % exempli gratia
\newcommand{\cf}{\textit{cf.}~} % confer
\newcommand{\apriori}{\textit{a priori}} % a priori
\newcommand{\viceversa}{\textit{vice versa}} % vice versa
\newcommand{\versus}{\textit{vs.}~} % versus
\newcommand{\aka}{\textit{a.k.a.}~} % also known as
\newcommand{\perse}{\textit{per se}} % per se
\definecolor{darkblue}{rgb}{0,0,0.7} % darkblue color
\definecolor{green}{RGB}{57,181,74} % green color
\definecolor{violet}{RGB}{147,39,143} % violet color
\newcommand{\red}{\color{red}} % red command
\newcommand{\blue}{\color{blue}} % blue command
\newcommand{\orange}{\color{orange}} % orange command
\newcommand{\green}{\color{green}} % green command
\newcommand{\darkblue}{\color{darkblue}} % darkblue command
\newcommand{\defn}[1]{\textsl{\darkblue #1}} % emphasis of a definition
\newcommand{\mathdefn}[1]{{\darkblue #1}} % emphasis of a definition
\newcommand{\yes}{\textcolor{green}{\ding{52}}}
\newcommand{\no}{\textcolor{red}{\ding{55}}}



%%%%%%%%%%%%%%%%%%%%%%%%%%%%%%%%%%%%%%


% Polytope
\newcommand{\pol}[1][P]{\mathsf{#1}}%polytope
\newcommand{\polQ}{{\pol[Q]}}
\newcommand{\polF}{{\pol[F]}}
\newcommand{\polC}{{\pol[C]}}
\newcommand{\polR}{{\pol[R]}}
\newcommand{\polS}{{\pol[S]}}
\newcommand{\polT}{{\pol[T]}}

\newcommand{\verts}{V}
\newcommand{\facets}{F}
\newcommand{\edges}{E}

\newcommand{\simplex}{\pol[\Delta]} %simplex
\newcommand{\cube}[1][n]{\pol[\square]_{#1}}
\newcommandx{\ZG}[1][1=G]{\pol[Z]_{#1}} %Graphical zonotope
\newcommand{\permuto}[1][n]{\pol[\Pi]_{#1}} %Permutahedron

\newcommand{\PH}[1][\H]{\pol_{#1}}

% Hypergraphs
\renewcommand{\H}{\c H}
\newcommand{\reducedH}[1][X]{\c H\left|_{#1}\right.}
\newcommand{\contractH}[1][X]{\c H/_{(V\ssm #1)}}


\newtotcounter{nbProjects}
\setcounter{nbProjects}{0}
\let\oldsubsection\subsection
\renewcommand{\subsection}{\stepcounter{nbProjects}\oldsubsection}

\newtotcounter{nbDone}
\setcounter{nbDone}{0}
\newcommand{\done}{\stepcounter{nbDone}}

\newtotcounter{nbWriting}
\setcounter{nbWriting}{-1}
\DeclareRobustCommand{\InWriting}{\refstepcounter{nbWriting}\textcolor{red}{$\bullet$}\,}


\title{Projects}
\author{Germain POULLOT}
\date{Dec. 2025 $\to$ \today}

\begin{document}

ERASE THIS PAGE AFTER READING!

\bigskip

This document is made for tracking what I do, what I've been asked about, etc.
Nothing is mandatory, everything is fluid, and each sub-section is more or less an idea that would make a paper if I manage to develop it (yet, two sub-sections may be merged into one paper, or one sub-section may lead to several papers).
I tried to organize these sub-sections into meaningful sections

Each sub-section comes with:
\begin{compactitem}
\item a title
\item some tags (the list of my possible tags and their meaning is on the beginning of next page)
\item a list of co-authors
\item maybe one idea that summarizes what the project is about
\end{compactitem}
Put the tags in the title, and it will automatically be counted: the total appears under the title of the document (the count on the next page is currently false because of the present page...).

Most likely, the important project have their own Overleaf/TexPage/LaTeX folder, so I don't need to write too much about it.
Besides, I'm very good at remembering what we've said on a topic, just from the title; and I don't care about re-do fifteen times the same reasoning, so forgetting is not a problem.

Yet, there are some stuffs that I never remember: the status of my submitted papers, so I'm very exhaustive about it.
My advice, if you re-use this template, is to write the stuffs that you will forget and the stuffs for which forgetting will be a real problem.
The less you write, the more up-to-date your file will be.

\hrule
\medskip

The following is are two typical examples of sub-sections:

\subsection*{$\star$ Paths on edge polytope of random graphs}
with Torben and Martina

cf Overleaf


\subsection*{$\circ$ Minkowski decomposability of Johnson solids}
alone

Acknowledgments:
Laura Casabella,
Zoe Geiselmann,
Frederic Chapoton,
Arnau Padrol


Number of rays (i.e. triangles) of the $\mathbb{DC}$ of regular $(2n+1)$-gon:
$\frac{n(n+1)(2n+1)}{6}$

$\dots$ of the regular $2n$-gon (segments \& triangles):
$n + \frac{2n(n-1)(n-2)}{6}$

Remember: re-try implementation with NumberFields


\hrule
\medskip

When I finish a project, it is ``done'', so I put it in the first section.
I write very thoroughly the tracking info of when it has been submitted where, etc.

Such a sub-section finishes by ``$\backslash$done'': it serve for the automatic counting.




\clearpage





\maketitle


$\bullet$: just write it!

\InWriting: currently in the writing

$\circ$: alone

$\star$: people are counting on me to do something

$\square$: to be decided exactly what/who goes into the project

$\diamond$: project of other people that I need to watch over a bit (I'm not author)

$-$: too hard

\sout{blah}: done


{\centerline{\textbf{Number of current projects}: \total{nbProjects}}}

{\centerline{Among which are done: \total{nbDone} ; versus currently in writing: \fpeval{\totvalue{nbWriting}/2}}

\tableofcontents

\newpage

\setcounter{section}{-1}
\section{Article done, submitted, etc.}

\subsection{$\star$ \sout{Hamiltonicity of graph of acyclic orientations of a graph}}
with Leonie M\"uhlherr

On ArXiv (03/09/2026)

$\to$ Submitted to SIDMA (SIAM Journal on Discrete Maths) 15/09/2026 $\to$ rejected 29/09/26

$\to$ \emph{Combinatorial Theory}?

$\to$ FPSAC ?

\done

\subsection{$\circ$ \sout{Vertices of the monotone path polytopes of hypersimplices}}
Alone

Rejected \emph{European Journal of Combinatorics} 

Submitted \emph{Electronic Journal of Combinatorics} (03/09/2026)
(submitted with the review of the previous submission to \euro JC).

($\to$ Discrete \& Combinatorial Geometry next?)

\done

\subsection{\sout{Unimodality of the number of paths per length on polytopes}}
with Martina

Accepted \emph{International Mathematics Research Notices}

\done



\newpage

\section{Hypergraphs, hypergraphic polytopes}

\subsection{\InWriting $\star$ Diameter of hypergraphic polytopes}
With Félix, ongoing

cf Overleaf

\subsection{$\circ$ Random hypergraphic polytopes}
alone

Erd\"os-R\'enyi style!

\subsection{Code a nice class}

Good luck


\newpage
\section{3D polytopes}

\subsection{$\circ$ Minkowski decomposability of Johnson solids}
alone

Acknowledgments:
Laura Casabella,
Zoe Geiselmann,
Frederic Chapoton,
Arnau Padrol


Number of rays (i.e. triangles) of the $\mathbb{DC}$ of regular $(2n+1)$-gon:
$\frac{n(n+1)(2n+1)}{6}$

$\dots$ of the regular $2n$-gon (segments \& triangles):
$n + \frac{2n(n-1)(n-2)}{6}$

Remember: re-try implementation with NumberFields

\subsection{$-$ Fibonacci number of paths for 3D polytopes}
alone? with Dan?

Stuck


\newpage
\section{Indecomposability}
\subsection{$-$ Smilansky conjecture 2}
with Arnau?

Find an indecomposable polytope with ``few'' vertices in regards of its number of facets, in dimension $\geq 5$

annoying...


\subsection{$\circ$ $-$ An edge-dependencies style proof that $\mathbb{DC}($Asso$)$ is simplicial}
with Arnau?

no clue


\newpage
\section{Submodular cone, deformation cones}
\subsection{$\star$ Smooth polygons with simplicial deformation cone}

with Veronica Calvo Cortes

Use Oda's decomposition theorem.

Short paper, easy.

\subsection{$\circ$ Face-closed classes are negligible inside $\mathbb{SC}_n$}

Be careful : property should only depend on normal fan.

Examples: simple, bounded degree exedence, 

Easy to prove that the opposite class has $\geq 2^{2^n}$ elements.

\subsection{$\circ$ Size of $\c E(\permuto, \permuto)$}

Between $3\cdot 2^{n+1} - 2$ and $2^{(n+1)\cdot (n+1)!}$, but I need better.

If size ``small'', \ie same (or smaller) size as $\scr L(\mathbb{SC}_n^2)$, then the bound on the number of total faces of $\mathbb{SC}_n$ (and hence of rays), is doubly-exponential, \ie far better.

hard ?

\newpage
\section{Random polytopes}
\subsection{\InWriting $\star$ Paths on edge polytope of random graphs}
with Torben and Martina

cf Overleaf

\end{document}
